\tzpanchor{1116}
\tzpentryheader{1116}{1000up}

\begin{tzpsnapshot}
\tzpsnaprow{TZPID}{\tzpcode{ID1116}}
\tzpsnaprow{UUID}{\tzpcode{039442d4-c4ea-560e-958b-d43da30100a5}}
\tzpsnaprow{Type}{Transfer Statement}
\tzpsnaprow{Scale}{transfer}
\tzpsnaprow{LOC Classification}{Working routing: QA641 manifolds and topology; QC174.17 mathematical physics}
\tzpsnaprow{arXiv Classification}{Primary: math-ph; Cross-list: gr-qc, math.DG}
\end{tzpsnapshot}

\tzpsection{Canonical Statement}
HTrawin[psi, t] = Htunnel + HER + HSC + Hcymatic + HTZP ; A?tunnel = -hbar\textasciicircum{}2/(2 m\_\{eff\}) nabla\textasciicircum{}2 + V\_\{barrier\}(r, theta, phi, t) ; A?ER = integral sqrt\{|g|\} R\textasciicircum{}\{munu\} T\_\{munu\} dSigma ; A?SC = sum (epsilon\_k - mu)(c\textasciicircum{}dagger\_\{kuparrow\} c\_\{kuparrow\} + c\textasciicircum{}dagger\_\{kdownarrow\} c\_\{kdownarrow\}) + Delta sum (c\textasciicircum{}dagger\_\{kuparrow\} c\textasciicircum{}dagger\_\{-kdownarrow\} + h.c.)

\[
t] = Ĥtunnel + ĤER + ĤSC + Ĥcymatic + ĤTZP
\]
\[
Ĥtunnel = -\hbar^2/(2 m_{eff})  abla^2 + V_{barrier}(r, \theta, \phi, t)
\]

\noindent
\begin{minipage}[t]{0.48\linewidth}
\tzpsection{Wolfram Form}
\begingroup
\small
\[
\begin{aligned}
\mathfrak{W}_{\mathrm{TZP}}\!\left(\mathcal{K}_{1116}\right)
&\longmapsto
\left\langle \Omega^{\mathbb{T}},\; \Psi_{\mathrm{T}},\; \rho_{\mathrm{vac}} \right\rangle
\end{aligned}
\]
\[
\mathcal{K}_{1116}
:=
\operatorname{HoldForm}\!\left(\mathcal{T}_{1116}\right)
\]
\[
\begin{aligned}
\operatorname{Admissible}_{\mathrm{TZP}}\!\left(\mathcal{K}_{1116}\right)
&\Longleftrightarrow
\mathfrak{W}_{\mathrm{TZP}}\!\left(\mathcal{K}_{1116}\right)\not\equiv\varnothing
\end{aligned}
\]
\textbf{Wolfram register:} canonical symbol lift\\
\textbf{Title lift:} 1000up\\
\textbf{Symbol key:} \(\Omega^{\mathbb{T}}\): Trawinistic boundary curvature\\ \(\Psi_{\mathrm{T}}\): Trawinistic wavefunction\\ \(\rho_{\mathrm{vac}}\): vacuum energy density
\par
\endgroup
\end{minipage}\hfill
\begin{minipage}[t]{0.48\linewidth}
\tzpsection{TRAWIN Normalization}
\[
\tzpnormalizewrap{t] = Ĥtunnel + ĤER + ĤSC + Ĥcymatic + ĤTZP,\,Ĥtunnel = -\hbar^2/(2 m_{eff})  abla^2 + V_{barrier}(r, \theta, \phi, t),\,ĤER = \int \sqrt{|g|} R^{\mu u} T_{\mu u} d\Sigma}
\]

\small
\textbf{T}: Canonical Definition is localized within the signal arc of the directory.\\
\textbf{R}: Geometric Constraint preserves the operative relation that 1000up requires.\\
\textbf{A}: Aggregate Relation fixes the admissible closure condition for the entry.\\
\textbf{W}: 1000up is rendered as a reusable operator-level statement.\\
\textbf{I}: The informational content of 1000up is retained rather than flattened.\\
\textbf{N}: 1000up closes as a distinct normalized TRAWIN object.
\normalsize
\end{minipage}

\tzpsection{Derivable Consequences}
The entry now reads through actual sourced equations, so its local arc is carried by explicit mathematical relations rather than a uniform quaternary certificate record shell.
\clearpage

\tzpentryheader{1116}{Oxford Dictionary of the Queens, NY English}

\begingroup\small
\tzpsection{Definition}
1000up is a registry definition in Signal with secondary pressure from Engineering and Implementation Arcs.

% BEGIN TZP CONTEXTUAL MEANING
\tzpsection{Contextual Meaning}
\begingroup\footnotesize\setlength{\emergencystretch}{3em}\sloppy
\textbf{Formal statement:} HTrawin[psi, t] = Htunnel + HER + HSC + Hcymatic + HTZP ; A?tunnel = -hbar to the power 2/(2 m sub eff ) nabla to the power 2 + V sub barrier (r, theta, phi, t) ; A?ER = integral sqrt |g| R to the power munu T sub munu dSigma ; A?SC = sum (epsilon sub k - mu)(c to the power dagger sub kuparrow c sub kuparrow + c to the power dagger sub kdownarrow c sub kdownarrow ) + Delta sum (c to the power dagger sub kuparrow c to the power dagger sub -kdownarrow + h.c.). \textbf{Contextual meaning:} The entry reads 1000up as a controlled technical headword whose equation layer, proof route, and registry role are interpreted together. \textbf{Derivable consequences:} The entry now reads through actual sourced equations, so its local arc is carried by explicit mathematical relations rather than a uniform quaternary certificate record shell. \textbf{Lemma support:} Canonical Definition; Geometric Constraint; Aggregate Relation.\par
\endgroup
% END TZP CONTEXTUAL MEANING

\tzpsection{Technical Interpretation}
Technically, 1000up is read through the local equation chain and the TRAWIN normalization recorded on the entry page.

\tzpsection{Equation Grounding}
Equation anchors: t] =  tunnel +  ER +  SC +  cymatic +  TZP;   tunnel = -hbar to the power 2/(2 m sub eff ) abla to the power 2 + V sub barrier (r, theta, phi, t); and   ER = integral sqrt |g| R to the power mu u T sub mu u dSigma. Proof route: show that the stated equations form a coherent progression for the entry and remain compatible under the TRAWIN ordering. Formalization route: prove that the final closure relation follows from the preceding entry-specific equations.

\tzpsection{Interpretive Role}
The entry now reads through actual sourced equations, so its local arc is carried by explicit mathematical relations rather than a uniform quaternary certificate record shell.
\endgroup

\tzpsection{TZP Thesaurus}
\begin{enumerate}[leftmargin=1.6em,itemsep=6pt,topsep=4pt]
\item \textbf{1000 Up}: entry-specific alternate wording for indexing, related literature, and local formal context.
\item \textbf{1000 Up Formal Analogue}: entry-specific alternate wording for indexing, related literature, and local formal context.
\item \textbf{Up Equivalence}: entry-specific alternate wording for indexing, related literature, and local formal context.
\end{enumerate}

\tzpsection{Encyclopedia Mathematica}
\begin{samepage}
\noindent\begin{minipage}[t]{0.45\linewidth}
\vspace{0pt}
\raggedright
\scriptsize\textbf{Image reading.} The rendered manifold at right is the per-ID light plate for 1000up. It should be read as the visual companion to the entry's equation layer, not as a repeated template diagram.\par
\end{minipage}\hfill
\begin{minipage}[t]{0.53\linewidth}
\vspace*{-0.10in}
\raggedright
\tzpencyclopediaimage{1116}
\end{minipage}
\end{samepage}

\clearpage

\tzpentryheader{1116}{Direct Equation Progression and Proof Trajectory}

\tzpsection{Direct Equation Progression}
\begin{enumerate}[leftmargin=1.6em,itemsep=9pt,topsep=6pt]
\item \textbf{Canonical Definition}
\[
t] = Ĥtunnel + ĤER + ĤSC + Ĥcymatic + ĤTZP
\]
This fixes the first explicit relation carried by the entry rather than a generic certificate record normalization.

\item \textbf{Geometric Constraint}
\[
Ĥtunnel = -\hbar^2/(2 m_{eff})  abla^2 + V_{barrier}(r, \theta, \phi, t)
\]
This adds the next operational equation that advances the entry beyond its opening definition.

\item \textbf{Aggregate Relation}
\[
ĤER = \int \sqrt{|g|} R^{\mu u} T_{\mu u} d\Sigma
\]
This closes the progression with an actual admissibility or closure relation instead of recycling the normalization shell.
\end{enumerate}

\tzpsection{Proof}
\textbf{Proof route.} show that the stated equations form a coherent progression for the entry and remain compatible under the TRAWIN ordering.

\textbf{Formal method.} prove that the final closure relation follows from the preceding entry-specific equations.

\medskip
\tzpproofartifacts{Isabelle audit: corpus classification recorded in appendix.}{Lean audit: compiled corpus certificate.}{Rocq audit: compiled corpus certificate.}

\tzpsection{Dependencies and Test Handle}
\begin{tzpsnapshot}
\tzpsnaprow{Dependencies}{\tzpidref{1115}, \tzpidref{1100}}
\tzpsnaprow{Node}{\tzpcode{registry_id_1116}}
\tzpsnaprow{Support Titles}{Canonical Definition; Geometric Constraint; Aggregate Relation}
\tzpsnaprow{Test Handle}{If the cited entry equations cannot be made mutually admissible in one TRAWIN-normalized frame, the entry fails.}
\end{tzpsnapshot}

\tzpsection{Canonical Equation Links}
\tzpequationlinks
  {t] = Ĥtunnel + ĤER + ĤSC + Ĥcymatic + ĤTZP}
  {Ĥtunnel = \cdots}
  {ĤER = \int \sqrt{|g|} R^{\mu u} T_{\mu u} d\Sigma}
